\documentclass[12pt]{article}
\usepackage[utf8]{inputenc}
\usepackage[T1]{fontenc}
\usepackage{hyperref}
\title{Pelende Language Resources for K-12 Curriculum Development}
\author{Compiled by Claude Sonnet 5.5 (claude-sonnet-5-5) for Andr\'as Kornai}
\date{}
\begin{document}
\maketitle

\section{Introduction}
Pelende (ISO 639-3 ppp) is a Bantu (Niger-Congo) language of the Kenge area, Bandundu region, western Democratic Republic of the Congo. Glottolog places it in a Yaka--Pelende--Lonzo group; Wikipedia's Yaka article calls Pelende and Lonzo political rather than ethnolinguistic entities, so its status as a separate language is unresolved. Very little was found.

\section{Primary Educational Resources}
No primers, textbooks or school programmes specific to Pelende were found.

\section{Linguistic and Reference Materials}
\begin{itemize}
\item Glottolog entry for Pelende: \url{https://glottolog.org/resource/languoid/id/pele1244.bigmap.html}
\item Glottolog reference record, de Beir 1975: \url{https://glottolog.org/resource/reference/id/77157}
\item Wikipedia on Yaka (Congo--Angola), covering the Pelende label: \url{https://www.wikipedia.com/wiki/ISO_639:ppp}
\end{itemize}

\section{Religious and Community Resources}
\begin{itemize}
\item Joshua Project entry for Pelende (reports Scripture only as portions): \url{https://joshuaproject.net/people_groups/14353/CG}
\item Joshua Project language page for ppp: \url{https://joshuaproject.net/index.php/languages/ppp}
\end{itemize}

\section{Digital and Audio Resources}
No Pelende-specific audio, video or software was found.

\section{Curriculum Development Context}
Congolese schooling in this region typically relies on French and larger languages; no Pelende school materials are known from these searches. Related Yaka and Kikongo-cluster materials might be adaptable, but that would need confirmation with speakers.

\end{document}